% Imported from LaTeX_and_pdfs/DGDT_dump.tex; see provenance.json.
\part{Classical Mechanics applications}


cf. Arnold, Kozlov, Neishtadt (2006) \cite{dgdt:AKN2006}. 

If known forces $\mathbf{F}_1 \dots \mathbf{F}_n$ acts on points, then 
\begin{equation}
\sum_{i=1}^n \langle m_i \ddot{ \mathbf{r}}_i - \mathbf{F}_i , \mathbf{\xi}_i \rangle = 0
\end{equation}
cf. Eq. (1.26) of Arnold, Kozlov, Neishtadt (2006) \cite{dgdt:AKN2006}, where $\xi_1 , \dots \xi_n$ are arbitrary tangent vectors to $M$, $\mathbf{\xi}_i, \dots \mathbf{\xi}_n \in TM$.  

$\sum_{i=1}^n \langle m_i \ddot{ \mathbf{r}}_i - \mathbf{F}_i , \mathbf{\xi}_i \rangle$ called "general equation of dynamics" or d'Alembert-Lagrange principle.  

\section{Orbital Mechanics}

References:
Arnold, Kozlov, Neishtadt (2006) \cite{dgdt:AKN2006}
Choquet-Bruhat (2009) \cite{dgdt:YChoquet-Bruhat2009}, pp. 39, 3. General relativity
Bate, Mueller, White (1971) \cite{dgdt:BMW1971}, Sec. 1.2, pp. 5, \url{https://en.wikipedia.org/wiki/N-body_problem}
Abraham and Marsden (2008) \cite{dgdt:AbMa2008}, Ch. 9
Pr\'{a}staro (1996) \cite{dgdt:Pras1996}, Remark 3.20 many-bodies dynamics, pp. 346


\subsection{Conic sections}

The following is an amalgamation or summary of the development in Ch. 13, Sections 13.18-22, from pp. 497, Apostol (1991) \cite{dgdt:Apos1991}

\begin{definition}[General definition of conic sections (vector-formulated definition)]\label{dgdt:Def:ConicSectionsVectorFormulation}
Given line $L$, called the \emph{directrix}, pt. $F$, the \emph{focus}, such that $F \notin L$., \\
\phantom{Given} eccentric $e > 0$, \\
\phantom{Given} $d = $ distance of $L$ from $F$, then \\
the \textbf{conic section} $C = \lbrace X \rbrace$ (set of all points) such that 
\begin{equation}\label{dgdt:Eq:ConicSectionVectorGeneralDefinition}
\| X - F \| = ed(X, L)
\end{equation}
Now, \\
$e < 1 $ ellipse, \\
$e = 1$ parabola, \\
$e > 1$ hyperbola.

Further, if $N$ unit normal to $L$, then $\forall \, P \in L$, 
\begin{equation}
d(X, L) = | (X- P) \cdot N |
\end{equation}
If $(X- P) \cdot N \gtrless 0$, $X \in $ positive half plane (negative half-plane). \\
\end{definition} 

\begin{theorem}[Further conic section development]\label{dgdt:Thm:MoreConicSectionsFromVectorFormulation}
	Building upon Def. \ref{dgdt:Def:ConicSectionsVectorFormulation}, \\
	
	If $F \in $ positive half plane determined by $N$, then $F -d N \in L$.  \\
	
	If $F \in $ negative half plane determined by $N$, then $F + dN \in L$. \\
	
Choose this convenient, specific point: $ F + dN = P \in L$ \quad \, $\Longrightarrow \| X - F \| = e |(X - (F + dN))\cdot N |$ \\

If $F = 0$, $\| X \| = e| X\cdot N - d | = r = e(d - r\cos{\theta})$. Then
\begin{equation}\label{dgdt:Eq:ConicSectionEqWithCenterAtFocus}
\frac{ed}{r} = 1 + e \cos{\theta} \text{ or } \boxed{ r = \frac{ed }{ 1 + e\cos{\theta}} } \text{ for $F = 0$ }
\end{equation}	

Consider symmetry about the origin (i.e. of $X, \exists \, -X$). 

As a caveat, let $A = e(d + F\cdot N)$. Then
\[
\begin{gathered}
	\| X - F \| = e | (X - (F + dN )) \cdot N | = e d(X, L) = e | X \cdot N - \frac{A}{e} | 
\end{gathered}
\]
Then
\begin{equation}\label{dgdt:Eq:ConicSectionDistanceFromFocusSquared}
\| X - F \|^2 = \| X \|^2 - 2 X\cdot F + \| F \|^2 = e^2 ((X\cdot N)^2 - 2 \frac{A}{e} X\cdot N + \left( \frac{A}{e}\right)^2 )
\end{equation}
Using the symmetry about the origin, $X \to -X$, \\
\[
\begin{gathered}
	\Longrightarrow X\cdot F = Ae X\cdot N \text{ or } X \cdot (F - Ae N) = 0
\end{gathered}
\]
Then $F = AeN$. Note that $A = e(F\cdot N \pm d)$ (for $F$ in negative (positive) half-plane). 

Then
\[
\begin{gathered}
Ae = e^2 ( Ae + d) \text{ or } \frac{A}{e} = Ae + d \text{ or } A \left( \frac{1}{e} - e \right) = d \text{ or } A = \frac{ed}{1 - e^2 } \\
\Longrightarrow \| X - Ae N \| = e| (X - (AeN + dN)) \cdot N | = e|X\cdot N - (Ae + d) |
\end{gathered}
\]
Let $X = aN$. Plug this into Eq. \ref{dgdt:Eq:ConicSectionDistanceFromFocusSquared}.
\begin{equation}\label{dgdt:Eq:SemiMajorAxis}
\begin{gathered} 
a^2 - 2aA e + A^2 e^2 = e^2 (a^2 - \frac{2A}{e} a + \left( \frac{A}{e} \right)^2 ) = e^2 a^2 - 2Aae + A^2 \Longrightarrow a = A \\
\Longrightarrow a = \frac{ed}{1 - e^2 }
\end{gathered} 
\end{equation}

Thus, we've shown that $A$ is equal to the \emph{semi-major axis} $a$ and that 
\begin{equation}
\boxed{ F = aeN }
\end{equation}
Thus, the semi-major axis is obtained when $X = \pm aN$ (vertices), so that $\| X \|^2 + (ae)^2 = e^2 ( X\cdot N)^2 + a^2$ is satisfied. \\
Likewise, the semi-minor axis is obtained if $X = \pm bN_{\perp}$, so that
\begin{equation}\label{dgdt:Eq:ConicSectionSemiMinorAxis}
\begin{gathered}
	b^2 + (ae)^2 = 0 + a^2 \Longrightarrow \boxed{ b = a \sqrt{ 1 - e^2 } }
\end{gathered} 
\end{equation}

\end{theorem}

We can then rewrite Eq. \ref{dgdt:Eq:ConicSectionEqWithCenterAtFocus}, with Eq. \ref{dgdt:Eq:SemiMajorAxis} as
\[
\boxed{ r = \frac{ a(1-e^2) }{ 1 + e\cos{\theta} } }
\]

With Eq. \ref{dgdt:Eq:ConicSectionDistanceFromFocusSquared}, and using the symmetry of $X$ about the origin, then
\[
\| X \|^2 + \| F \|^2 = e^2 (X\cdot N)^2 + a^2
\]

Choose $\mathbf{N} = \mathbf{e}_x$. With $\|X\|^2 = x^2 + y^2$, and $\|F\| = Ae = ae$, then
\begin{equation}\label{dgdt:Eq:ByOriginSymmetryConicSection}
\begin{gathered}
x^2 + y^2 +a^2 e^2 = e^2 x^2 + a^2 \text{ or } x^2 ( 1 -e^2) + y^2 = a^2 ( 1-e^2) \\
\end{gathered}
\end{equation}

cf. 13.23 "Cartesian equations for the conic sections", Apostol (1991) \cite{dgdt:Apos1991}.

\subsubsection{Ellipse}

If $e <1$, then from Eq. \ref{dgdt:Eq:ByOriginSymmetryConicSection},

\[
\begin{gathered}
\xrightarrow{ \text{ if } e < 1 } \frac{x^2 }{ a^2 } + \frac{y^2}{b^2} = 1
\end{gathered}
\]

Observe that what  Greiner (2004) \cite{dgdt:Grei2004} in pp. 249, Ch. "Planetary Motions", (a) Ellipse, denotes $c$ is what we and Apostol denote as $\|F \| =ae$, the distance from the origin to the focus. Observe that
\[
c^2 + b^2 = a^2 e^2 + a^2 (1-e^2) = a^2
\]

\subsubsection{Hyperbola}

If $e > 1$, then from Eq. \ref{dgdt:Eq:ByOriginSymmetryConicSection},

\[
\begin{gathered}
\xrightarrow{ \text{ if } e > 1 } \frac{x^2 }{ a^2 } - \frac{y^2}{b^2} = 1
\end{gathered}
\]
Let $b = |a| \sqrt{ e^2 - 1}$. The foci of a hyperbola are at $(c,0)$, $(-c,0)$ where $c= |a| e = \sqrt{ a^2 + b^2}$. 

\subsubsection{Parabola}

cf. pp. 507 of  13.23 "Cartesian equations for the conic sections", Apostol (1991) \cite{dgdt:Apos1991}

Go back to Eq. \ref{dgdt:Eq:ConicSectionVectorGeneralDefinition}, where
\[
\| X - F \| = d(X, L)
\]
Let $L$ be a vertical line at $x = \mp c$, place focus $F = (\pm c, 0)$. \\
If $X = (x,y)$, $X-F = (x \mp c, y)$, so 
\[
\begin{gathered}
\| X- F \|^2 = (x \mp c)^2 + y^2 = | x \pm c|^2 = x^2 \pm 2xc + c^2 = x^2 \mp 2xc + c^2 + y^2
\Longrightarrow y^2 = \pm 4xc
\end{gathered}
\]
So midway between the focus and directrix, at the origin, is the vertex of the parabola.

\subsection{Central Force Problem via Energy Methods}

\subsubsection{Reduction to the Equivalent One-Body Problem}

cf. pp. 70 of Goldstein, Poole, Safko (2001) \cite{dgdt:GPS2001}

Let $\mathbf{R} = $ radius vector to the center of mass. \\

Let $\mathbf{r} = \mathbf{r}_2 - \mathbf{r}_1$ for the relative position of two masses. Then the Lagrangian
\begin{equation}
\mathcal{L} = T(\dot{\mathbf{R}}, \dot{ \mathbf{r}} ) - U (\mathbf{r}, \dot{\mathbf{r}}, \dots )
\end{equation}
cf. Eq. (3.1) of Goldstein, Poole, Safko (2001) \cite{dgdt:GPS2001}

The kinetic energy $T$ is the sum of the kinetic energy of the center of mass motion plus the kinetic energy of motion, about the center of mass $T'$.
\[
T = \frac{1}{2} (m_1 + m_2) \dot{ \mathbf{R}}^2 + T' \text{ with } T' = \frac{1}{2} m_1 \dot{ \mathbf{r}_1'}^2 + \frac{1}{2} m_2 \dot{\mathbf{r}_2'}^2
\]

$\mathbf{r}_1', \mathbf{r}_2'$ are radii vectors relative to the center of mass, since \\
\[
\begin{gathered}
\mathbf{R} = \frac{m_1 \mathbf{r}_! + m_2 \mathbf{r}_2 }{M}
\end{gathered}
\]
so 
\begin{equation}
\begin{aligned}
& \mathbf{r}_1 - \mathbf{R} = \frac{m_2 (\mathbf{r}_1 - \mathbf{r}_2 )}{ M} = -\frac{m_2}{M} \mathbf{r} \equiv \mathbf{r}_1' \\
&  \mathbf{r}_2' \equiv \mathbf{r}_2 - \mathbf{R} = \frac{m_1(\mathbf{r}_2 - \mathbf{r}_1 )}{ M} = \frac{m_1}{M} \mathbf{r}
\end{aligned}
\end{equation}
cf. Eqn. (3.2) of Goldstein, Poole, Safko (2001) \cite{dgdt:GPS2001}.

Thus $T' = \frac{1}{2} \frac{m_1 m_2}{m_1 + m_2 } \dot{ \mathbf{r}}^2$

So
\begin{equation}
\mathcal{L} = \frac{M}{2} \dot{\mathbf{R}}^2 + \frac{1}{2} \frac{m_1 m_2}{M} \dot{\mathbf{r}}^2 - U (\mathbf{r}, \dot{\mathbf{r}}, \dots )
\end{equation}
cf. Eqn. (3.3) of Goldstein, Poole, Safko (2001) \cite{dgdt:GPS2001}.

3 coordinates $\mathbf{R} $ are cyclic, so CM is at rest or moving uniformly. Drop the first term. The rest of the Lagrangian is exactly what would be expected if we had a fixed center of force with a single particle at distance $\mathbf{r}$ from it.

Let 
\[
\mu = \frac{m_1 m_2}{ m _1 + m_2} \equiv \text{ reduced mass }
\]

\subsubsection{Equations of Motion from Lagrangian and First Integrals}

cf. Sec. 3.2 "The Equations of Motions and First Integrals, Chapter 3 The Central Force Problem, pp. 72 of Goldstein, Poole, Safko (2001) \cite{dgdt:GPS2001}.

Goldstein, Poole, Safko argues: assume potential is of form $V(r)$ and so $V(r)$ has only radial distance, so the problem has spherical symmetry. Then, Goldstein, et. al. argues $\mathbf{L} = \mathbf{r} \times \mathbf{p}$ conserved. Thus $\mathbf{r}$ and $\mathbf{p} \perp \mathbf{L}$ always. 

So 
\begin{equation}
\mathcal{L} = T- V = \frac{1}{2} \mu (\dot{r}^2 + r^2 \dot{\theta}^2 ) - V(r)
\end{equation}
cf. Eqn. (3.6) Goldstein, Poole, Safko (2001) \cite{dgdt:GPS2001}.

Now
\[
p_{\theta} = \frac{ \partial \mathcal{L} }{ \partial \dot{\theta}} = \mu r^2 \dot{\theta}
\]
Since $\frac{d}{dt} \frac{\partial \mathcal{L}}{ \partial \dot{\theta} } = \frac{\partial \mathcal{L}}{\partial \theta}$, then
\begin{equation}
\dot{p}_{\theta} = \frac{d}{dt} ( \mu r^2 \dot{\theta} ) = 0 = \frac{\partial \mathcal{L}}{\partial \theta}
\end{equation}
cf. Eqn. (3.7)  Goldstein, Poole, Safko (2001) \cite{dgdt:GPS2001}.

So 
\begin{equation}
\mu r^2 \dot{\theta} = l
\end{equation}
cf. Eqn. (3.8)  Goldstein, Poole, Safko (2001) \cite{dgdt:GPS2001}, where $l$ is the constant magnitude of the angular momentum.

So then $\frac{d}{dt} \left( \frac{1}{2} r^2 \dot{\theta} \right) = 0$ where $\frac{1}{2} r^2 \dot{\theta}$ is the areal velocity, the area swept out by radius vector per unit time.

Goldstein argues that $\frac{1}{2} r^2 \dot{\theta}$ follows from $dA = \frac{1}{2} r(rd\theta)$, so $\frac{dA}{dt} = \frac{1}{2} r^2 \frac{d\theta}{dt}$. Thus, angular momentum conservation is equal to areal velocity constant. Do keep in mind that this argument is based upon the diagram in  Goldstein, Poole, Safko (2001) \cite{dgdt:GPS2001} and presumably the small angle approximation (perhaps using $\sin{\theta} \approx \theta$ for small $\theta$). Thus, Kepler's 2nd. law: the radius vector sweeps out equal areas in equal times.

For the $r$ coordinate:

\[
\begin{gathered}
\frac{d}{dt} \left( \frac{ \partial \mathcal{L} }{ \partial \dot{r}} \right) = \frac{\partial \mathcal{L}}{\partial r} \Longrightarrow \frac{d}{dt} ( \mu \dot{r} ) = \mu r\dot{\theta}^2 - \frac{\partial V}{\partial r} \text{ or } \mu \ddot{r} - \mu r \dot{\theta}^2 = -\frac{\partial V}{\partial r}
\end{gathered}
\]

Denote 
\begin{equation}
f(r) = \mu \ddot{r} - \mu r\dot{\theta}^2
\end{equation}
cf. Eqn. (3.11) of  Goldstein, Poole, Safko (2001) \cite{dgdt:GPS2001}.

Since $\mu r^2 \dot{\theta} = l$, then
\begin{equation}
f(r) = \mu \ddot{r} - \frac{l^2}{\mu r^3}
\end{equation}
cf. Eqn. (3.12) of Goldstein, Poole, Safko (2001) \cite{dgdt:GPS2001}.

Since 
\begin{equation}
\mu \ddot{r} = \frac{- \partial}{\partial r} \left( V + \frac{1}{2} \frac{l^2 }{\mu r^2} \right)
\end{equation}
cf. Eqn. (3.14) of Goldstein, Poole, Safko (2001) \cite{dgdt:GPS2001} and since $\mu \ddot{r} \dot{r} = \frac{d}{dt} \left( \frac{1}{2} \mu \dot{r}^2 \right)$ and $\frac{d}{dt} g(r) = \frac{\partial g}{\partial r} \dot{r}$, then
\[
\frac{d}{dt} \left( \frac{1}{2} \mu \dot{r}^2 \right) = \frac{-d}{dt} \left( V + \frac{1}{2} \frac{l^2 }{ \mu r^2 } \right) \text{ or } \frac{d}{dt} \left( \frac{1}{2} \mu \dot{r}^2 + \frac{1}{2} \frac{l^2}{ \mu r^2 } + V \right) = 0
\]
and so
\begin{equation}
\frac{1}{2} \mu \dot{r}^2 + \frac{1}{2} \frac{l^2}{ \mu r^2} + V = \text{ constant } 
\end{equation}
cf. Eqn. (3.15) of  Goldstein, Poole, Safko (2001) \cite{dgdt:GPS2001}, with
\[
\frac{1}{2} \frac{l^2 }{\mu r^2 } = \frac{1}{2 \mu r^2} \mu^3 r^4 \dot{\theta}^2 = \frac{\mu r^2 \dot{\theta}^2 }{2} 
\]
and so
\begin{equation}
\frac{1}{2} \mu \dot{r}^2 + \frac{1}{2} \mu r^2 \dot{\theta}^2 + V = E = T+V 
\end{equation}
cf. Eqn. (3.13) of  Goldstein, Poole, Safko (2001) \cite{dgdt:GPS2001}.

For $E = \frac{1}{@} \mu \dot{r}^2 + \frac{l^2}{2\mu r^2 } + V$, then 

\begin{equation}
\dot{r} = \sqrt{ \frac{2}{\mu} \left( E - V - \frac{l^2}{2 \mu r^2}  \right)  }
\end{equation}
cf. Eqn. (3.16)  of Goldstein, Poole, Safko (2001) \cite{dgdt:GPS2001}, or 

\begin{equation}
dt =\frac{dr}{ \sqrt{ \frac{2}{\mu} \left( E - V - \frac{l^2}{2 \mu r^2 } \right) } }
\end{equation}
cf. Eqn. (3.17) of  Goldstein, Poole, Safko (2001) \cite{dgdt:GPS2001}.

Let
\[
t =\int_{r_0}^r \frac{dr}{ \sqrt{ \frac{2}{\mu} \left( E - V - \frac{l^2}{2 \mu r^2 } \right) } }
\]
with initial value $r_0$.

From $l= \mu r^2 \dot{\theta}$,
\begin{equation}
d\theta = \frac{l dt}{\mu r^2} 
\end{equation}
cf. Eqn. (3.19) of Goldstein, Poole, Safko (2001) \cite{dgdt:GPS2001} so for $\theta_0 \equiv $ initial value of $\theta$.

\begin{equation}
\theta = l \int_0^t \frac{dt'}{ \mu r^2(t') } + \theta_0
\end{equation}
cf. Eqn. (3.20) of Goldstein, Poole, Safko (2001) \cite{dgdt:GPS2001}.

In Sec. 3.3, "The Equivalent 1-Dimensional Problem and Classification of Orbits" of  Goldstein, Poole, Safko (2001) \cite{dgdt:GPS2001}, consider $V ' = V + \frac{1}{2} \frac{l^2}{\mu r^2}$, so that $F' = \frac{-\partial V'}{\partial r} = f(r) + \frac{l^2}{\mu r^3}$. \\

Let $V = \frac{-k}{r}$, e.g. $V = \frac{-GMm}{r} $ with $k = GMm$. \\

Then consider $V' = \frac{-k}{r} + \frac{l^2}{2\mu r^2}$ and its plot.


Compare this, Thm. \ref{dgdt:Thm:MoreConicSectionsFromVectorFormulation}, against (26.6) of Greiner (2004) \cite{dgdt:Grei2004}:
\begin{equation}\label{dgdt:Eq:ConicSectionPolarCoordinatesFocusCenter}
r = \frac{ed}{ 1 + e\cos{\theta}}, \quad \, e = \frac{C\widehat{L}_0^2}{GM}, \quad \, d = 1/C
\end{equation} cf. Eq. (26.32) of Greiner (2004) \cite{dgdt:Grei2004}. 

Investigate which physical quantities (e.g. $E, L$) eccentricity $e$ depends on.

\[
\begin{gathered}
u = \frac{GM}{ \widehat{L_0}^2 } + C \cos{(\theta - \theta_0)} \Longrightarrow \frac{du}{d\theta} = -C \sin{(\theta - \theta_0)} \\
\Longrightarrow \begin{gathered}
\frac{1}{2} m L_0^2 \left[ C^2 \sin^2{ (\theta - \theta_0)} + \left( \frac{GM}{L_0^2} \right)^2 + 2 C \frac{GM}{L_0^2} \cos{(\theta - \theta_0)} + C^2 \cos^2{(\theta - \theta_0)} \right] = \\
= \frac{1}{2} m \widehat{L}_0^2 \left[ C^2 + \left( \frac{GM}{\widehat{L}_0^2 } \right)^2 + 2 C \frac{GM}{ \widehat{L}_0^2 } \cos{(\theta - \theta_0)} \right] = \\
= E - V\left( \frac{1}{u} \right) = E + GMmu = E + \frac{(GM)^2 }{ \widehat{L}_0^2 } m + C GM m \cos{(\theta - \theta_0)}
\end{gathered} \\
\Longrightarrow \frac{1}{2} m \widehat{L}_0^2 C^2 = E + \frac{1}{2} \frac{ (GM)^2 }{ \widehat{L_0}^2 } m 
\end{gathered} 
\]

Solve for the dummy constant $C$:
\begin{equation}
C = \sqrt{ \frac{2E}{m \widehat{L}_0^2 } + \frac{ (GM)^2 }{ \widehat{L}_0^4 } }
\end{equation}
cf. Eq. (26.36) of Greiner (2004) \cite{dgdt:Grei2004}.

From Eq. (26.32) of Greiner (2004) \cite{dgdt:Grei2004} or Eq. \ref{dgdt:Eq:ConicSectionPolarCoordinatesFocusCenter}, 
\[
\begin{gathered}
e = \frac{C \widehat{L}_0^2 }{ GM} = \sqrt{ \frac{2 E }{ m \widehat{L}_0^2} + \frac{ (GM)^2 }{\widehat{L}_0^4 } } \left( \frac{\widehat{L}_0^2 }{ GM} \right)  
\end{gathered}
\]
So
\begin{equation}
\boxed{ e =  \sqrt{ \frac{ 2 E \widehat{L}_0^2 }{ m (GM)^2 } + 1 } }
\end{equation}
cf. Eq. (26.37) of Greiner (2004) \cite{dgdt:Grei2004}. 

Hence, shape of the path depends on total energy $E$, angular momentum $L_0 = m \widehat{L}_0$ of moving body (recall that $\widehat{L} := \mathbf{r} \times \mathbf{v}$). \\

It holds that for a

\[
\begin{aligned} 
\text{ parabola } & e = 1, & E = 0 & \\  
\text{ ellipse } & 0 < e < 1, & E < 0, & \qquad \, \frac{ (GM)^2 m }{ -2 \widehat{L}_0^2 } < E < 0 \\  
\text{ circle } &  e = 0 & E = \frac{ -(GM)^2 m }{ 2\widehat{L}_0^2 } & \\
\text{ hyperbola } & e = 1, & E > 0 & \, 
\end{aligned} 
\]

The general form of initial energy $E_0$ is, recall, give by
\[
E_0 = \frac{1}{2} m v^2 + V(r) = \frac{1}{2} mv^2 - \frac{GMm}{r_0} = m \left( \frac{v^2}{2} - \frac{GM}{r_0} \right) 
\]

Notice that for the case of the parabola that \\
$E_0 =0 $ when 
\begin{equation}\label{dgdt:Eq:EscapeVelocityAsParabola}
v = \sqrt{ \frac{2GM}{r_0} } 
\end{equation}
Eq. \ref{dgdt:Eq:EscapeVelocityAsParabola} is the expression for the \textbf{escape velocity}.

Knowing what we know now, let's rewrite some of the expressions from before:
\[
\begin{gathered} 
r(\theta) = \frac{ \widehat{L}_0^2 / GM}{ 1 + \frac{C \widehat{L}_0^2 }{GM} \cos{(\theta - \theta_0)} } \\ 
\widehat{L}_0^2 := r_0^2 \dot{\theta}_0 \\ 
C = \sqrt{\frac{ 2 E}{ m \widehat{L}_0^2 } + \frac{ (GM)^2 }{ \widehat{L}_0^4 } }
\end{gathered} 
\]

\subsection{Path parameters (i.e. orbital parameters)}

cf. pp. 262, Greiner (2004) \cite{dgdt:Grei2004}.

\[
a = \frac{1}{2} ( r(\theta = 0) + r(\theta = \pi) ) = \frac{1}{2} \left[ \frac{ \widehat{L}_0^2 / GM }{ 1 + e } + \frac{\widehat{L}_0^2 / GM }{ 1- e} \right] = \frac{ \widehat{L}_0^2 / GM}{ 1 - e^2 }
\]
The addition above for $a$ is needed because we center origin at a focus, for this expression. 

\begin{equation}\label{dgdt:Eq:Semi-MajorAxis}
a = \frac{\widehat{L}_0^2 /GM}{ \frac{ - 2 E \widehat{L}_0^2 }{ m (GM)^2 } } = \frac{ GMm }{ - 2 E} 
\end{equation}
cf. Eq. (26.39) of Greiner (2004) \cite{dgdt:Grei2004}. 

\begin{equation}
b = \sqrt{ a^2 - f^2 } = \sqrt{a^2 - e^2 a^2 } = a \sqrt{ 1 - e^2 } = a \sqrt{ \frac{ -2 E \widehat{L}_0^2 }{ m (GM)^2 } } = \widehat{L}_0 \sqrt{ \frac{ - m}{ 2E} }
\end{equation}
cf. Eq. (26.40) of Greiner (2004) \cite{dgdt:Grei2004}. 

We can also calculate the perigee (point of closest approach to mass $M$):
\begin{equation}
r(\theta =0 ) = \frac{ \widehat{L}_0^2 / GM }{ 1 + e } = \frac{ \widehat{L}_0^2 / GM }{ 1 + \sqrt{ \frac{ 2 E \widehat{L}_0^2 }{ m (GM)^2 } + 1 } }
\end{equation}

Given Eq. \ref{dgdt:Eq:Semi-MajorAxis} for $a$, plug in for $E$:
\begin{equation}
\begin{gathered}
a = \frac{ - GMm}{2 E} = - \frac{GMm}{ 2 \left( \frac{1}{2} m v_0^2 - \frac{GMm}{r} \right) } = \frac{ -GM }{ v_0^2 - \frac{2GM}{r} } = \frac{ - GM }{ v_0^2 \left( 1 - \frac{2 GM}{ rv^2 }\right)} 
\end{gathered}
\end{equation}
Notice when the denominator is 0. It is when $v_0$ is the the escape velocity (cf. Eq. \ref{dgdt:Eq:EscapeVelocityAsParabola}). 

\subsubsection{Eccentricity from $r_0$, $\mathbf{v}_0$}

Consider being given $\mathbf{r}$, $\mathbf{v}$, with $\mathbf{r}$ in a coordinate frame with its origin at the center of mass (or, if $M$ large enough, from the focus where $M$ is centered at). 

The first step is to use a coordinate \emph{frame} (set of orthonormal unit vectors) with $\mathbf{r}$, $\mathbf{v}$ forming a \emph{plane}. Note this important fact from geometry: a plane is spanned by 2 non-collinear vectors. So $\mathbf{r}, \mathbf{v}$ better not be parallel to each other.

Let 
\[
\begin{aligned} 
	& \mathbf{v} = v^x \mathbf{e}_x + v^y \mathbf{e}_y + v^z \mathbf{e}_z = v^r \mathbf{e}_r + v^{\varphi} \mathbf{e}_{\varphi} \\ 
	& \mathbf{r} = x \mathbf{e}_x + y \mathbf{e}_y + z \mathbf{e}_z = r\mathbf{e}_r
\end{aligned} 
\]
Now 
\[
\begin{gathered}
(\mathbf{r} \times \mathbf{v})_i = \epsilon_{jki} r_j v_k 
\end{gathered}
\]
To get the vector \emph{orthogonal} to the plane that is formed by $\mathbf{r}, \mathbf{v}$, 
\[
\begin{gathered}
\begin{gathered}
	((\mathbf{r} \times \mathbf{v}) \times \mathbf{r})_m = \epsilon_{ijm} \epsilon_{jki} r_j v_k r_l = \epsilon_{ilm} \epsilon_{ijk} r_j v_k r_l = \\
	= r^2 v_m - v_k r_k r_m = r^2 v_m - v_k r_k r_m
\end{gathered}
\text{ or } ((\mathbf{r} \times \mathbf{v}) \times \mathbf{r}) = r^2 \mathbf{v} - (\mathbf{v} \cdot \mathbf{r}) \mathbf{r}
\end{gathered}
\]

Then
\[
\begin{gathered}
\widehat{L}_0 = r^2 \dot{\theta} = r \left( \mathbf{v} \cdot  \frac{ ((\mathbf{r} \times \mathbf{v}) \times \mathbf{r}) }{ | (\mathbf{r} \times \mathbf{v} ) \times \mathbf{r} | } \right) = r \left( \frac{r^2 v^2 - (\mathbf{v} \cdot \mathbf{r})^2 }{ |r^2 \mathbf{v} - (\mathbf{v} \cdot \mathbf{r}) \mathbf{r} | } \right) 
\end{gathered}
\]

Thus, we can find the following:
\[
\begin{gathered} 
e^2 = 1 + \frac{ 2 E \widehat{L}_0^2 }{ m (GM)^2 } = 1 + \frac{ 2 \left( \frac{1}{2} m v_0^2 - \frac{GM m }{r} \right) \widehat{L}_0^2 }{ m (GM)^2 }  = 1 + \frac{\left( v_0^2 - \frac{2GMm}{r} \right) \widehat{L}_0^2 }{ (GM)^2 }
\end{gathered} 
\]

\[
\begin{gathered} 
a = \frac{-GMm}{2E} = -\frac{GMm}{ 2 \left( \frac{1}{2} mv_0^2 - \frac{GMm}{r} \right) } = \frac{-GM}{v_0^2 - \frac{2GM}{r} } = \frac{-GM}{ v_0^2 \left( 1 -  \frac{2GM}{rv_0^2 } \right) }
\end{gathered} 
\]


TODO: Calculate period of revolution. 

Here are some calculations (forms) that can help with further calculations:
\begin{equation}
\begin{gathered} 
	\widehat{L}_0^2 = b^2 \left( \frac{2E}{-m} \right) \text{ or } \frac{\widehat{L}_0^2}{GM} = b^2 \left( \frac{2E}{-GMm } \right) = \frac{b^2}{a} = a(1-e^2)
\end{gathered} 
\end{equation}
using Eq. \ref{dgdt:Eq:ConicSectionSemiMinorAxis}, for $e^2 = 1 - \left( \frac{b}{a} \right)^2$. 


\subsection{Path in coordinates} 

For $\mathbf{r}$ centered at focus $F$ and $\mathbf{r}, \mathbf{v}$ expressed in coordinates in the plane of the conic section path, i.e. in the plane formed by $\mathbf{r}, \mathbf{v}$, with the axes expressed as $\frac{\partial}{\partial x}$, $\frac{ \partial }{ \partial y}$, 
\begin{equation}
\mathbf{r} = r\cos{\theta} \frac{\partial }{ \partial x} + r\sin{\theta} \frac{\partial}{\partial y}
\end{equation}
Compare this to Eq. (2.5-1) of Sec. 2.5.1 "Expressing $r$ and $v$ in the Perifocal System" of Bate, Mueller, White (1971) \cite{dgdt:BMW1971}.

Recall the equation for the conic section in polar coordinates, with the coordinates centered at focus $F$:
\begin{equation}
r = \frac{a(1-e^2) }{ 1 + e\cos{\theta}}
\end{equation}
Compare this to Eq. (1.5-4) of Bate, Mueller, White (1971) \cite{dgdt:BMW1971}.

Now
\begin{equation}
\mathbf{v} = \dot{\mathbf{r}} = (\dot{r} \cos{\theta} - r\dot{\theta} \sin{\theta}) \frac{\partial}{\partial x} + (\dot{r}\sin{\theta} + r\dot{\theta}\cos{\theta}) \frac{\partial }{ \partial y}
\end{equation}
Compare this to the equation on pp. 72, of Sec. 2.5.1 "Expressing $r$ and $v$ in the Perifocal System" of Bate, Mueller, White (1971) \cite{dgdt:BMW1971}.

Now
\[
\begin{gathered}
\dot{r} = \frac{ a ( 1 - e^2) }{ (1+ e\cos{\theta})^2 } e\sin{\theta} \dot{\theta} = \frac{r^2}{a (1-e^2) } e\sin{\theta} \frac{ \widehat{L}_0 }{ r^2} = \frac{ \widehat{L}_0 e\sin{\theta}}{ a(1-e^2)} = \frac{\widehat{L}_0 e\sin{\theta} }{ \frac{ -GMm}{2E} \left( \frac{ L_0^2}{ \left( \frac{(GM)^2 m }{ -2E} \right) } \right) } = 
\end{gathered}
\]
\begin{equation}
\Longrightarrow \dot{r} = \frac{GM}{\widehat{L}_0} e\sin{\theta}
\end{equation}
Compare this to Eq. (2.5-2) of Sec. 2.5.1 "Expressing $r$ and $v$ in the Perifocal System" of Bate, Mueller, White (1971) \cite{dgdt:BMW1971}.

Then
\[
\begin{gathered}
	\Longrightarrow \frac{ a (1-e^2) }{ (1 + e\cos{\theta})^2 } \dot{\theta} = \frac{GM}{ \widehat{L}_0 } \text{ or } 
\end{gathered}
\]
\begin{equation}
r\dot{\theta} = \frac{GM}{\widehat{L}_0 } (1 + e\cos{\theta})
\end{equation}
Compare this to Eq. (2.5-3) on pp. 73 of Sec. 2.5.1 "Expressing $r$ and $v$ in the Perifocal System" of Bate, Mueller, White (1971) \cite{dgdt:BMW1971}.

Then
\[
\begin{gathered}
\mathbf{v} = \frac{GM}{\widehat{L}_0} \left[ \left( e \sin{\theta} \cos{\theta} - (1 + e\cos{\theta}) \sin{\theta} \right) \frac{\partial }{ \partial x} + \left( e\sin^2{\theta} + (1 + e\cos{\theta} ) \cos{\theta} \right) \frac{\partial }{ \partial y} \right] = 
\end{gathered}
\]
\begin{equation}
\mathbf{v} = \frac{GM}{ \widehat{L}_0} \left[ \sin{\theta} \frac{\partial}{\partial x} + (e + \cos{\theta}) \frac{\partial }{\partial y} \right] = \frac{GM}{ \widehat{L}_0} \left[ \frac{y}{r} \frac{\partial}{\partial x} + (e + \frac{x}{r}) \frac{\partial }{\partial y} \right]
\end{equation}
Compare this to Eq. (2.5-4) on pp. 73 of Sec. 2.5.1 "Expressing $r$ and $v$ in the Perifocal System" of Bate, Mueller, White (1971) \cite{dgdt:BMW1971}.

Consider this ansatz from subsection 4.4.3 "The $f$ and $g$ expressions from Ch. 4 section on "Position and velocity - A function of time" in Bate, Mueller, White (1971) \cite{dgdt:BMW1971}. 

Because $\mathbf{r}, \mathbf{v}$ lie in the plane (i.e. coplanar) containing the entire conic section trajectory, then make the following ansatz:
\begin{equation}
\begin{aligned}
& \mathbf{r} = f\mathbf{r}_0 + g \mathbf{v}_0 \\ 
& \dot{\mathbf{r}} =: \mathbf{v} = \dot{f} \mathbf{r}_0 + \dot{g} \mathbf{v}_0
\end{aligned} 
\end{equation}

Solve for $f$ and $g$ by taking the appropriate cross products:
\begin{equation}\label{dgdt:Eq:fAndgSeparatedByCrossProducts}
\begin{aligned} 
	& \mathbf{r} \times \mathbf{v}_0 = f\mathbf{r}_0 \times \mathbf{v}_0 \\ 
	& \mathbf{r} \times \mathbf{r}_0 = g \mathbf{v}_0 \times \mathbf{r} 
\end{aligned}
\end{equation}

\subsection{Eccentricity anomaly; Eccentricity anomaly coordinates}

Let the point of interest be P. Choose the origin for the coordinates $(x, y)$ to have \emph{its origin} at the \emph{center of the ellipse}. 

Let point C be the point at the center of the ellipse (origin), let point A be the projection of point P onto the $x$-axis, point F be the point at the focus.

Recall that the length from the center C to the focus, $\overline{\text{CF}}$, is $ae$, where $a$ is the semi-major axis, and eccentricity is $e$. 

The eccentricity anomaly $E$ (it's just an angle) is given by
\begin{equation}
\cos{E} = \frac{x}{a}
\end{equation}
Remember $E$ is \emph{not} the angle formed by triangle CAP, but rather, first take the line through P and A. This line is perpendicular to the x-axis. Then intersect that line with a circle, centered at C, of radius $a$. Then, the line from C to that intersection forms an angle $E$ with the $x$-axis.

By using the ellipse equation and trigonometry,
\[
\begin{gathered}
	\frac{x^2}{a^2} + \frac{y^2}{b^2} = 1 \Longrightarrow \cos^2{E}  + \frac{y^2}{b^2} = 1 \text{ or } \frac{y^2}{b^2} = \sin^2{E} \text{ or } 
\end{gathered}
\]

\begin{equation}
\sin{E} = \frac{y}{b}
\end{equation}

Now
\[
\overline{\text{AF}} = \overline{\text{CF}} - \overline{\text{CA}} = ae - a\cos{E}
\]
Use the right triangle PAF and the hypotenuse being $r$ (distance from the focus F to point P) to find
\[
\overline{\text{AF}} = r\cos{(\pi - \theta)} = -r\cos{\theta}
\]
So
\[
\cos{\theta} = \frac{ a ( \cos{E} - e) }{r} 
\]
Then, plug into the conic section equation, Eq. \ref{dgdt:Eq:ConicSectionEqWithCenterAtFocus}, 
\[
\begin{gathered}
	r = \frac{ed}{ 1 + e\cos{\theta}} = \frac{ a (1-e^2) }{ 1 + e\cos{\theta}} = \frac{b^2 /a}{ 1 + e\cos{\theta}} = \frac{b^2/a}{ 1 + e\left( \frac{ a(\cos{E} - e) }{ r} \right) } \\
		\Longrightarrow 1 = \frac{ b^2 / a}{ r + e(a(\cos{E} - e))} \\
		\Longrightarrow r = \frac{b^2}{a} - e(a(\cos{E} - e)) = \frac{b^2}{a} + ae^2 - ea \cos{E} = a(1 - e \cos{E})
\end{gathered}
\]
\begin{equation}\label{dgdt:Eq:EllipseEquationWithEccentricAnomaly}
\boxed{ r = a(1 - e\cos{E}) }
\end{equation}
Compare this to Eq. (4.6-9) on pp. 215 of Sec. 4.6 in "The Kepler Problem" of Bate, Mueller, White (1971) \cite{dgdt:BMW1971}.

For the coordinates with the origin at the \emph{focus} $F$, it is given by 
\begin{equation}\label{dgdt:Eq:EllipticalPathInEccentricAnomaly}
\begin{aligned}
& x = a(\cos{E} - e) \\
& y = b\sin{E} = a\sqrt{ 1 - e^2} \sin{E}
\end{aligned}
\end{equation}
Compare this to Eqns. (4.6-7), (4.6-8) on pp. 215 of Sec. 4.6 in "The Kepler Problem" of Bate, Mueller, White (1971) \cite{dgdt:BMW1971}.

Indeed, for $E=0$, $x = a - ae$ (the distance from the focus F to the vertex is $a - ae$), $E = \pi/2$, $x =-ae$ (the distance from the focus F and back to the center C) and for $E = \pi$, $x = -a - ae$ (the distance to the other vertex).

\subsubsection{Path in terms of eccentric anomaly}

Take the time derivative of Eq. \ref{dgdt:Eq:EllipticalPathInEccentricAnomaly}:
\begin{equation}
\begin{aligned}
& \dot{x} = -a \sin{E} \dot{E} \\ 
& \dot{y} = a\sqrt{1 - e^2} \cos{E} \dot{E} 
\end{aligned} 
\end{equation}
To determine $\dot{E}$, consider both expressions for the radial distance $r$ for Eq. \ref{dgdt:Eq:EllipseEquationWithEccentricAnomaly} and Eq. \ref{dgdt:Eq:ConicSectionEqWithCenterAtFocus}:

\[
\begin{gathered} 
\begin{gathered} 
	r = \frac{a (1-e^2) }{ 1 + e\cos{\theta}} = \frac{b}{ 1 + e\cos{\theta}} \\
	\Longrightarrow \dot{r} = \frac{ -b}{ (1+e\cos{\theta})^2} (-e\sin{\theta})\dot{\theta} = \frac{ be \sin{\theta} \dot{\theta} }{ (1+ e\cos{\theta})^2 } = \\
		\dot{r} = \frac{r^2 e}{b} \sin{\theta} \dot{\theta}
\end{gathered} \\ 
\begin{gathered}
r = a(1 - e\cos{E}) \\
\Longrightarrow \dot{r}  = +ae\sin{E} \dot{E} 
\end{gathered} \\
\begin{gathered}
\Longrightarrow ae \sin{E} \dot{E} = \frac{r^2 e}{b} \sin{\theta} \dot{\theta} \text{ or } a\sin{E} \dot{E} = \frac{r^2}{b} \sin{\theta} \dot{\theta} = \frac{ \widehat{L}_0 \sin{\theta} }{ b} = \frac{\widehat{L}_0 }{b} \frac{y}{r} \\
\Longrightarrow \dot{E} = \frac{ \widehat{L}_0 }{ab} \frac{y}{r\sin{E}} = \frac{\widehat{L}_0 } {ar} = \frac{ b}{ a^{3/2} r} 
\end{gathered}
\end{gathered} 
\]
Compare the expression for $\dot{E}$ with Eq. (4.6-11) on pp. 216 of Bate, Mueller, White (1971) \cite{dgdt:BMW1971}.

Armed with an expression for $\dot{E}$, plug this back into the expression for $\dot{r}$:
\begin{equation}
\dot{r} = a e \sin{E} \left( \frac{ \widehat{L}_0 }{ ar} \right) = \frac{ e\sin{E} \widehat{L}_0 }{ r}
\end{equation}
Compare this to the equation on pp. 216, "Position and Velocity - a function of time" in Ch. 4 of Bate, Mueller, White (1971) \cite{dgdt:BMW1971}. 

Use the expressions above in Eqns. \ref{dgdt:Eq:fAndgSeparatedByCrossProducts}. For $\mathbf{r}, \mathbf{r}_0$ centered at focus F:
\[
\begin{gathered} 
xv_{0, y} - y v_{0,x} = a(\cos{E} - e) b \cos{E_0} \dot{E}_0 - b\sin{E} (-a\sin{E_0} \dot{E}_0) = \\
	= ab\dot{E}_0 \left[ ( \cos{E} - e) \cos{E_0} + \sin{E} \sin{E_0} \right] = \frac{b^2}{ \sqrt{a} r} \left[ \cos{(E - E_0)} - e\cos{E_0} \right] = \frac{ \widehat{L}^2 \left( \frac{-m}{2E} \right) }{ \sqrt{GM} \sqrt{ \frac{m}{-2E} } r } \left[ \cos{(E- E_0)} - e\cos{E_0} \right] = \\
	 = \frac{\widehat{L}_0 b }{ r \sqrt{ GM} } \left[ \cos{ (E- E_0)} - e\cos{E_0} \right] 
\end{gathered} 
\]
Comparing this against the "right-hand side (RHS)" with $f (\mathbf{r}_0 \times \mathbf{v}_0)$, the RHS has a magnitude of $f \widehat{L}_0$ in the same direction as $(\mathbf{r} \times \mathbf{v}_0)$. Thus
\begin{equation}
f = \frac{ b }{ r \sqrt{ GM} } \left[ \cos{ (E- E_0)} - e\cos{E_0} \right]
\end{equation}

For the expression involving $g$ in Eqns. \ref{dgdt:Eq:fAndgSeparatedByCrossProducts}, 
\[
\begin{gathered}
	\mathbf{r} \times \mathbf{r}_0 = xy_0 - yx_0 = a(\cos{E} - e) b \sin{E_0} - b\sin{E} a (\cos{E_0} - e) = -ab \left[ (\cos{E_0} ) \sin{E} - (\cos{E}-e) \sin{E_0} \right] 
\end{gathered}
\]
with
\[
\begin{gathered}
ab = a \widehat{L}_0 \sqrt{ \frac{-m}{2E} } = a\widehat{L}_0 \left( \frac{a}{GM} \right)^{1/2} = \left( \frac{a^3}{GM} \right)^{1/2} \widehat{L}_0
\end{gathered}
\]
For the RHS with $g \mathbf{v}_0 \times \mathbf{r}_0$, $\mathbf{v}_0 \times \mathbf{r}_0$, by the right-hand rule or the right-hand orientation, it is in the \emph{opposite} direction as $\mathbf{r}\times \mathbf{r}_0$. Thus, apply a negative side on the RHS.

\begin{equation}
\Longrightarrow g = \left( \frac{ a^3}{GM} \right)^{1/2} \left[ (\cos{E_0} - e) \sin{E} - (\cos{E} - e) \sin{E_0} \right] 
\end{equation}

